% !PRINTED 638
\weberpage{638}
By the theorem of \gesperrt{Lagrange} (§.155) every function
$\Theta_\nu$ possessing the four properties $\alpha)$, $\beta)$,
$\gamma)$, $\delta)$ can be expressed rationally, in the field $\Omega$
of the metacyclic functions of $x$, by one of them, $\omega_\nu$,
provided only that the $\omega_0, \omega_1 \ldots \omega_{n-2}$ are
different from one another; one may therefore choose for these functions
the $f_\nu$.

For if $u$ denotes a variable, and if

\weq{19}{(u - \omega_0)(u - \omega_1) \cdots (u - \omega_{n-2}) = \varphi(u),}

then the function $\varphi(u)$ admits the permutations $s, t$ and also
the substitution $\sigma = (\varepsilon, \varepsilon^g)$; and hence, by
Theorem 3. of §.183, it is an integral function of $u$ with coefficients
metacyclic in regard to $x$, and independent of $\varepsilon$.

If now $\Theta_0, \Theta_1 \ldots \Theta_{n-2}$ is a system of
functions possessing the properties $\alpha) \beta) \gamma) \delta)$,
then the sum

\weq{20}{\sum_{0,\, n-2}^{\nu} \frac{\Theta_\nu\, \varphi(u)}{u - \omega_\nu} = \chi(u)}

is an integral function of degree $(n-2)$ of $u$, which likewise admits
the substitutions $s, t, \sigma$ and which therefore has coefficients in
$\Omega$. If one sets

\weq{21}{\frac{\chi(u)}{\varphi'(u)} = \Theta(u)}

and then sets $u = \omega_\nu$ in (20), it follows that

\[
\Theta_\nu = \Theta(\omega_\nu),
\]

where $\Theta$ denotes a rational function whose coefficients are
metacyclic functions of the $x$, and no longer contain $\varepsilon$.

\begin{center}
§. 185.

\gesperrt{Roots of metacyclic equations.}
\end{center}

In the formulae of the preceding section, the roots $\xi_0, \xi_1
\ldots \xi_{n-1}$ of a metacyclic equation of degree $n$, irreducible
in some field $\mathfrak{K}$, are now to be introduced in place of the
variables $x_0, x_1 \ldots x_{n-1}$, and indeed in such an order that
the metacyclic functions of the $\xi$ are rational (in $\mathfrak{K}$),
which by the theorems of the preceding section is always possible. But
in doing so we first make two further restrictive as-
%CONT>>
